% =====================================================================
%  Section 6 -- Discussion
% =====================================================================

\section{Discussion}
\label{sec:discussion}

\subsection{What the negative KD result means}
\label{sec:discussion-negative-kd}

On FSSD, knowledge distillation does not help when the student is an ImageNet-pretrained lightweight UNet. We frame this as a \emph{negative result} rather than a \emph{null result} because the comparison is not ``we did not find a positive effect'' but rather ``the best student (MobileViT-XXS, no-KD) \emph{exceeds} the teacher (UNet$+$EfficientNet-B0)''. Distillation on this benchmark is at best redundant and at worst counter-productive: combined KD on MobileViT-XXS produced $-0.005$ IoU vs.~the same architecture trained without it.

We interpret this mechanistically. The classical motivation for KD, Hinton's \emph{dark knowledge} argument~\cite{hinton2015distilling} and subsequent feature-map alignment work~\cite{romero2015fitnets,liu2019structured,wang2020intra,he2021channel} assumes that the student has insufficient capacity or pretraining to recover the teacher's behaviour from labelled data alone, so the teacher's soft probabilities and intermediate features provide a richer supervision signal than the binary mask. On FSSD, that assumption is violated in two directions: (i)~the benchmark has only 663 images and saturates at IoU $\approx 0.93$, so the labelled signal is already nearly sufficient; and (ii)~the student's ImageNet pretraining already supplies the visual priors (edge detectors, texture filters, semantic-region features) that a teacher trained on the same FSSD would have been transferring. The teacher's ``dark knowledge'' is therefore largely redundant with the student's pretraining, and the response-based KL term reduces to a noise-injection regulariser. The feature-MSE term comes closest to the no-KD baseline (MobileViT-XXS: $-0.0011$ IoU, against $-0.0041$ for response-only and $-0.0049$ for combined KD; Table~\ref{tab:t3-kd-ablation}), plausibly because it carries at least \emph{some} signal beyond the labels, but it still does not beat the baseline.

This finding has two implications for the flood-segmentation literature:

\begin{enumerate}
\item \textbf{Benchmarks matter.} A KD recipe that wins on Cityscapes (19 classes, 5000 images, IoU $\approx 0.80$--$0.85$) does not transfer to FSSD (2 classes, 663 images, IoU $\approx 0.93$). Future KD work in this space should report headroom (teacher--best-pretrained-student gap) \emph{before} claiming distillation effects.
\item \textbf{Modern ImageNet pretraining is doing the work.} Current \texttt{timm}-pretrained lightweight backbones brings the no-KD baseline so high that compression strategies \emph{other than} distillation (quantization, pruning, architecture choice) dominate the deployment-relevant trade-offs. We recommend that authors investing in flood-segmentation compression budget their effort accordingly.
\end{enumerate}

The combined-corpus experiment (Section~\ref{sec:results-combined}) supports the headroom reading and adds a caveat. Where the teacher out-segments the no-KD students on flood pixels (0.781 vs.\ 0.695--0.754 flood-only IoU), combined KD closes most of the gap, with the gain significant for all three students. It also transmits the teacher's false alarms: distilled students fire on the teacher's false-alarm scenes 2.3--2.6 times as often as their no-KD counterparts, and the increase is specific to those scenes. A student trained to imitate the teacher imitates its errors too. A teacher should therefore be audited on the failure mode that matters for deployment, here false alarms on dry scenes, before it is used for distillation, not only on the headline IoU that selects it.

\subsection{Per-platform deployment recommendations}
\label{sec:discussion-deployment}

Based on the joint segmentation $\times$ latency $\times$ stability evidence (Sections~\ref{sec:results-kd}--\ref{sec:results-latency}), we make the following deployment recommendations:

\begin{itemize}
\item \textbf{Field-laptop / tablet inference (Apple Silicon class):} EfficientNet-Lite0 UNet INT8. 17\,ms p50 (\deployableMTwoFPS~FPS) at IoU 0.885, $3.83\times$ size reduction. Stable across folds ($\pm 0.015$) even under the per-tensor recipe. This is the configuration we ship in our public release.
\item \textbf{Field-laptop / desktop inference (x86 class):} FP32, not INT8. On the Ryzen host INT8 is $2.8$--$3.0\times$ slower than FP32 for every network (Table~\ref{tab:t13-host2}), so the fastest configurations are FP32 MobileNetV3-Small (15\,ms) and FP32 EfficientNet-Lite0 (17\,ms); INT8 is worth its $3.4$--$3.8\times$ smaller file only when download size, not speed, is the constraint. Any INT8 deployment should be timed on the target CPU before it is shipped.
\item \textbf{Browser citizen-science / web-portal inference:} MobileNetV3-Small UNet FP32 (133\,ms p50; 7.5\,FPS on the M2 host), or MobileViT-XXS UNet FP32 (149\,ms p50; 6.7\,FPS) for slightly higher accuracy. Under ONNX Runtime Web~$\geq 1.18$ the INT8 artifacts load and run $1.3$--$1.4\times$ faster than FP32 (Table~\ref{tab:t13-host2}); per-channel INT8 MobileNetV3-Small (IoU 0.904, 3.8\,MB) is then the fastest and smallest browser option.
\item \textbf{Maximum-accuracy off-line inference:} MobileViT-XXS UNet FP32 (28\,ms p50 on M2~Pro CPU; IoU 0.937). Recommended for human-in-the-loop annotation pipelines.
\end{itemize}

The case against deploying the teacher (UNet$+$EffB0) is weaker than a latency-only comparison suggests. On a common ONNX Runtime CPU backend the teacher runs at 37.6\,ms p50 (\teacherMTwoFPS~FPS) on the M2~Pro and 23.4\,ms on the Ryzen, only $1.1$--$1.9\times$ slower than the FP32 students, and with per-channel scales it INT8-quantizes to IoU $0.917 \pm 0.005$ at 6.3\,MB (Table~\ref{tab:t12-ptq-mechanism}), which is more accurate than any INT8 student. What remains of the case is latency on the hardware where INT8 pays: the recommended EfficientNet-Lite0 INT8 is $2.2\times$ faster than the FP32 teacher on the M2~Pro (17.3 vs 37.6\,ms) while recovering 95\% of its IoU, and we did not time the INT8 teacher on that host. Where INT8 does not pay (the x86 host), the FP32 students are $1.1$--$1.5\times$ faster than the teacher, and a practitioner who can afford 23\,ms per image loses nothing by deploying the teacher itself.

Table~\ref{tab:t7-decision} consolidates these recommendations into a decision matrix, pairing each operational scenario with the configuration we recommend and the metric that decides it.

All of these recommendations carry the zero-shot caveats of Sections~\ref{sec:results-ood} and~\ref{sec:results-zeroshot}: the artifacts are validated for in-domain, close-range RGB; FSSD-trained models false-alarm on the majority of water-free scenes and must be retrained on negative-image-bearing data before \emph{unattended} monitoring use (Section~\ref{sec:results-combined} shows that this retraining brings student false-alarm rates down to 3--9\%, and that no-KD students are the safer choice there); and INT8 artifacts in particular should be re-validated on in-deployment-domain imagery, since their degradation worsens under distribution shift in a way that in-domain calibration metrics do not predict. The INT8 recommendations above also assume a CPU on which INT8 is faster than FP32; Table~\ref{tab:t13-host2} shows one on which it is not.

% =====================================================================
%  Table T7 -- Deployment Decision Matrix
%  Source of truth: summary.json (FP32 IoU), quantized.json (INT8 IoU),
%  lat_m2.json / lat_wasm.json (latency), T1 (size), T6 (OOD).
%  Re-presentation of the Section 6.2 recommendations; no new experiments.
% =====================================================================
\begin{table*}[ht]
\centering
\caption{\textbf{Deployment decision matrix.} Recommended configuration per operational scenario, with the deciding metrics. All values are drawn from Tables~\ref{tab:t1-footprints}, \ref{tab:t2-segmentation}, \ref{tab:t4-quant}, \ref{tab:t5-latency}, \ref{tab:t6-ood}, and~\ref{tab:t10-combined}; this table re-presents the Section~\ref{sec:discussion-deployment} recommendations as an at-a-glance practitioner reference and introduces no new experiments. Latency is p50 on the platform relevant to each scenario (Apple~M2~Pro ORT-CPU, or browser WASM SIMD on the M2 host); Table~\ref{tab:t13-host2} gives the x86 counterparts.}
\label{tab:t7-decision}
\small
\setlength{\tabcolsep}{4pt}
\begin{adjustbox}{max width=\textwidth}
\begin{tabular}{p{2.9cm}lcccp{6.2cm}}
\toprule
\textbf{Deployment scenario} & \textbf{Recommended} & \textbf{IoU} & \textbf{Latency (p50)} & \textbf{Size (MB)} & \textbf{Rationale / caveat} \\
\midrule
Field laptop / tablet (Apple-Silicon CPU)
  & \textbf{EffLite0-UNet INT8} & 0.885 & 17\,ms (M2) & 5.2
  & Fastest operating point on Apple Silicon (INT8 $2\times$ faster than FP32); the only student stable under per-tensor INT8 ($\pm 0.015$). On the x86 host INT8 is $2.8\times$ \emph{slower}: use FP32 there (Table~\ref{tab:t13-host2}). FSSD-trained, it false-alarms on 93\% of dry scenes (Table~\ref{tab:t9-zeroshot}): analyst-in-the-loop use only. \\
\addlinespace
Browser / citizen-science web portal
  & MNV3-S-UNet FP32 & 0.927 & 133\,ms (WASM) & 13.8
  & Fastest under single-threaded WASM SIMD; MobileViT-XXS FP32 (149\,ms, IoU 0.937) if accuracy matters more. Under ORT-Web~1.22 INT8 runs $1.3$--$1.4\times$ faster (Table~\ref{tab:t13-host2}); per-channel INT8 MobileNetV3-Small (IoU 0.904, 3.8\,MB) is then the smallest option. \\
\addlinespace
Maximum-accuracy off-line / human-in-the-loop
  & MobileViT-XXS-UNet FP32 & 0.937 & 28\,ms (M2) & 12.0
  & Highest IoU of any configuration; FP32 only (its INT8 loses $\geq 0.22$ IoU under every full-quantization recipe we tested, Table~\ref{tab:t12-ptq-mechanism}). \\
\addlinespace
Unattended monitoring (mostly dry scenes)
  & EffLite0-UNet FP32, no KD & 0.736\textsuperscript{$\ddagger$} & 34\,ms (M2) & 19.8
  & Must be retrained with water-free images: FSSD-trained models false-alarm on 70--100\% of dry scenes, combined-trained Lite0 on 3\% (Table~\ref{tab:t10-combined}). Lite0 over MobileNetV3-S because its INT8 variant is stable under either recipe (per-channel: FAR 0.051, flood IoU 0.625), whereas MobileNetV3-S is usable in INT8 only with per-channel scales (Table~\ref{tab:t11-combined-int8}). Skip KD: it inherits the teacher's false alarms. \\
\addlinespace
Cross-sensor / satellite optical
  & \emph{retrain / domain-adapt} & 0.16\textsuperscript{$\dagger$} & -- & --
  & No FSSD-trained model transfers zero-shot (Table~\ref{tab:t6-ood}); requires domain adaptation or in-sensor training data. \\
\bottomrule
\end{tabular}
\end{adjustbox}
\\[3pt]
{\footnotesize \textsuperscript{$\dagger$}Zero-shot Sen1Floods11 IoU; all models $\approx 0.16$ (Section~\ref{sec:results-ood}). \textsuperscript{$\ddagger$}Flood-only IoU on the FloodNet test split after combined-corpus retraining (Section~\ref{sec:results-combined}); not comparable with the FSSD IoUs above.}
\end{table*}


\subsection{The PTQ practitioner trap is the default recipe, not the architecture}
\label{sec:discussion-ptq-trap}

The 0.45-IoU PTQ collapse on MobileNetV3-Small and MobileViT-XXS would not be predicted from the classification-quantization literature, which typically reports about one point of top-1 accuracy lost to INT8 PTQ on ResNet/MobileNet/EfficientNet families on ImageNet~\cite{krishnamoorthi2018quantizing,jacob2018quantization}. That literature, however, assumes per-channel weight quantization for depthwise-separable networks, and Krishnamoorthi~\cite{krishnamoorthi2018quantizing} already reports that per-tensor weights drive MobileNet-family models to near-zero accuracy. Our ablation (Section~\ref{sec:results-quant}) reproduces exactly this on segmentation: the fold-dependent collapse of MobileNetV3-Small ($0.772$, $0.0001$ and $0.657$ IoU across folds) and the $0.19$--$0.30$ IoU loss of the teacher both disappear once the weight quantizer is per-channel, and keeping the HardSwish layers in FP32 adds nothing. We had attributed these collapses to the HardSwish activation range, which we had measured and found to be $12$--$17\times$ wider than ReLU6's; the measurement was correct and the inference from it was not. Two features of the segmentation setting still make the per-tensor failure worse than in classification: the output is a per-pixel binary mask, so a biased post-sigmoid threshold shifts the whole segmentation, and the collapse is fold-dependent, so a single validation fold can pass an artifact that fails on the next.

MobileViT-XXS is the genuine architectural case, but not for the reason the transformer-quantization literature would suggest. Its attention and LayerNorm blocks, which that literature identifies as the fragile components~\cite{lin2022fqvit,li2023repqvit}, quantize at no measurable cost here; the loss sits in the two high-resolution MobileNetV2-style stages before the first transformer block, which per-channel scales do not fix. The loss is not a property of the FSSD weights either: the same per-channel recipe collapses the combined-trained MobileViT-XXS outright in two of three folds (Table~\ref{tab:t11-combined-int8}). Why these stages fail in MobileViT-XXS and not in MobileNetV3-Small or the teacher, whose early blocks are of the same type, is open; the differences we can point to are MobileViT-XXS's very narrow early channel widths (16--24) and its SiLU activations in those blocks, and distinguishing them needs per-layer activation statistics we have not collected.

The practical guidance has therefore changed. A practitioner who runs the default ONNX Runtime static-quantization recipe on a lightweight segmentation backbone and sees it collapse should first switch \texttt{per\_channel} on; for every convolutional network in this study that one setting turns a collapse into a 0.01--0.04 IoU loss, on both training corpora and under zero-shot transfer (Tables~\ref{tab:t9-zeroshot} and~\ref{tab:t11-combined-int8}). Only if the loss persists, as it does for MobileViT-XXS, is the architecture the problem, and the next steps are keeping the identified layers in FP32 or quantization-aware training. In either case the INT8 artifact should be validated on every fold and on out-of-domain imagery (Section~\ref{sec:results-zeroshot}), and timed on the target CPU, since INT8 is not faster everywhere (Section~\ref{sec:results-latency}).

\subsection{Qualitative failure modes}
\label{sec:discussion-failures}

Figure~\ref{fig:f3-qualitative} shows that both teacher and EfficientNet-Lite0-INT8 student share two persistent failure modes:
\begin{itemize}
\item \textbf{Bright-reflective-surface confusion:} wet asphalt, glare-on-roof, and large reflective puddles are misclassified as flood. This is a labelling-modality issue, not a model-capacity issue: RGB alone cannot disambiguate water from wet non-water surfaces.
\item \textbf{Vegetation-occluded edges:} when flooded water is partially overhung by trees or tall grass, both models under-segment the flood boundary. The 5-pixel margin at the boundary contributes disproportionately to the residual IoU gap.
\end{itemize}

Both failure modes argue for SAR or multispectral fusion rather than for a different RGB segmentation recipe. We discuss this in Section~\ref{sec:limitations}.

Table~\ref{tab:t8-taxonomy} organises the failure modes observed across the whole study (the two RGB-modality failures above, the cross-sensor distribution shift of Section~\ref{sec:results-ood}, and the INT8 PTQ failures of Section~\ref{sec:results-quant}) into a single taxonomy, noting for each its mechanism, error direction, the models affected, and the mitigation it implies. The categories are qualitative: we report the error direction (which matters for the operational cost asymmetry of Section~\ref{sec:discussion-asymmetry}) but do not claim per-class frequencies.

% =====================================================================
%  Table T8 -- Error Taxonomy
%  Qualitative synthesis of failure modes established in
%  Sections 5.3 (PTQ), 5.7 (OOD), 6.4 (qualitative failure modes).
%  No new experiments; per-class frequencies are NOT claimed.
% =====================================================================
\begin{table*}[ht]
\centering
\caption{\textbf{Error taxonomy for FSSD flood segmentation.} Four failure classes observed across this study, their mechanism, the direction of the resulting error (which direction costs more depends on the use: false negatives in event-time mapping, false positives in unattended monitoring, Section~\ref{sec:discussion-asymmetry}), the models affected, and the mitigation each implies. Categories are qualitative, synthesised from Sections~\ref{sec:results-quant}, \ref{sec:results-ood}, and~\ref{sec:discussion-failures}; we do not claim per-class frequencies.}
\label{tab:t8-taxonomy}
\small
\setlength{\tabcolsep}{4pt}
\begin{adjustbox}{max width=\textwidth}
\begin{tabular}{p{2.6cm}p{4.2cm}p{2.1cm}p{2.4cm}p{3.2cm}}
\toprule
\textbf{Failure class} & \textbf{Mechanism} & \textbf{Error type} & \textbf{Affected models} & \textbf{Mitigation} \\
\midrule
Reflective-surface confusion
  & RGB cannot disambiguate water from wet non-water surfaces (wet asphalt, roof glare, standing puddles)
  & False positive
  & Teacher and all students
  & SAR / multispectral fusion (outside the RGB scope) \\
\addlinespace
Vegetation-occluded boundary
  & Flood edges overhung by trees or tall grass are under-segmented
  & False negative
  & Teacher and all students
  & Modality limitation; boundary-aware loss at best partial \\
\addlinespace
Cross-sensor distribution shift
  & Close-range RGB $\rightarrow$ Sentinel-2 satellite optical domain gap
  & Whole-image collapse (IoU\,$\approx$\,0.16)
  & All (uniformly)
  & Domain adaptation or in-sensor retraining \\
\addlinespace
INT8 PTQ instability
  & A single per-tensor weight scale cannot cover the per-channel weight ranges of these lightweight encoders; in MobileViT-XXS the early high-resolution convolutions also fail under per-channel scales (attention blocks quantize cleanly)
  & Model-level collapse ($\Delta$IoU $-0.22$ to $-0.92$, fold-dependent)
  & MobileNetV3-Small, MobileViT-XXS, teacher (per-tensor); MobileViT-XXS only (per-channel)
  & Per-channel weight scales (fixes every CNN); keep MobileViT-XXS's early convolutions FP32, or QAT \\
\bottomrule
\end{tabular}
\end{adjustbox}
\end{table*}


\subsection{Operational error asymmetry}
\label{sec:discussion-asymmetry}

The metrics we report, IoU and F1, weight false positives and false negatives equally, but in flood response the two errors carry very different operational costs. Which error costs more depends on the use. During an event, when an analyst maps a flood known to be under way, a false negative (unmapped flood) can translate into under-warning and under-allocation of rescue resources, whereas a false positive mainly costs excess caution; here the false negative is the more expensive error. In unattended monitoring of scenes that are dry most of the time, the balance reverses: each false alarm triggers a dispatch or a review, and repeated false alarms erode trust in the warning itself, so a model that flags water in 70--100\% of dry scenes (Section~\ref{sec:results-zeroshot}) is unusable however good its recall. We report precision and recall separately in Table~\ref{tab:t2-segmentation} precisely so this trade-off is visible rather than hidden inside a single symmetric score: the FP32 teacher is close to balanced (precision $0.962$, recall $0.960$), and the two qualitative failure modes above pull in opposite directions: bright-surface confusion inflates false positives, vegetation occlusion inflates false negatives.

INT8 conversion further perturbs this balance, and not in a stable direction: the deployable EfficientNet-Lite0 student's per-fold (precision, recall) after quantization are $(0.965, 0.934)$, $(0.917, 0.963)$, and $(0.897, 0.963)$, recall-favouring (safer, over-segmenting) in two folds but precision-favouring (missing flood) in the third. We therefore make two recommendations for event-time mapping. First, tune the post-sigmoid decision threshold toward recall, accepting more false positives to reduce missed flood, rather than leaving it at the default $0.5$; do this only for models trained with water-free images (Section~\ref{sec:results-combined}), since for FSSD-trained models it would worsen an already disqualifying false-alarm rate. Second, re-calibrate that threshold \emph{after} INT8 conversion, since quantization shifts the operating point by an amount and in a direction our results show is not predictable from the FP32 model alone. A full cost-weighted evaluation (e.g.\ an $F_\beta$ score with $\beta>1$, or a region-level warning-recall metric) is left to future work: the per-image cost model is application- and jurisdiction-specific and cannot be fixed once for all deployments.
